\section{Method}
\label{sec:method}
\subsection{Problem formulation}
For an engine at cycle $t$, input $\mathbf{X}_t\in\mathbb{R}^{W\times d}$ stacks the last $W{=}30$ cycles of $d{=}17$ channels (three operating settings and 14 sensors, min--max scaled with global statistics). The model $f_{\mathbf{w}}$ regresses the RUL $y_t=\min(R_{\max},\,T-t)$ with $R_{\max}=125$ (piece-wise linear target) and $T$ the engine's failure cycle.

\subsection{1D-CNN regressor}
The network (22\,273 parameters) is: Conv1d(17$\to$32, $k{=}3$)--ReLU--MaxPool(2); Conv1d(32$\to$64, $k{=}3$)--ReLU--MaxPool(2); flatten (448); Linear(448$\to$32)--ReLU; Linear(32$\to$1). The float32 model is 89\,092 bytes; one downlink plus one uplink is 178\,184 tensor bytes per round, excluding protocol overhead. These are architecture-derived payload sizes; the Pi profiler records measured process RSS and training time.

\subsection{Federated averaging}
With $K$ clients holding $n_k$ windows ($N=\sum_k n_k$), round $t$ proceeds: the server sends $\mathbf{w}^{(t)}$; each client runs $E$ local epochs of Adam on its data minimising the configured local objective $\mathcal{L}$ (MSE by default, or \texttt{AsymmetricRULLoss}), giving $\mathbf{w}_k^{(t+1)}$; the server aggregates
\begin{equation}
\mathbf{w}^{(t+1)}=\sum_{k=1}^{K}\frac{n_k}{N}\,\mathbf{w}_k^{(t+1)}.
\end{equation}
FedProx adds $\frac{\mu}{2}\lVert\mathbf{w}-\mathbf{w}^{(t)}\rVert_2^2$ to each client's local objective, with $\mu=0.01$. The FedAdam server uses the pseudo-gradient $\mathbf{g}_t=\mathbf{w}^{(t)}-\sum_k(n_k/N)\mathbf{w}_k^{(t+1)}$, Adam moments $(\beta_1,\beta_2)=(0.9,0.99)$, and server step size $0.01$.

\subsection{Paired-suite ablations}
The paired FD001--FD004 suite's score-aware ablation averages the same exponential penalty used by the NASA metric over each batch. To bound gradients, the exponent is capped at 20 during training; evaluation uses the uncapped NASA score. We label its centralized and federated variants ``Central + exp.'' and ``FedAvg + exp.'' in the generated tables. The CNN-LSTM ablation uses a 16-channel convolution, max pooling and a 24-unit LSTM head (4\,889 parameters); it is trained centrally against the 1D-CNN centralized baseline.

\subsection{Quadratic asymmetric risk-aware loss}
The separate controlled FD001/FD002 comparison uses \texttt{AsymmetricRULLoss} to avoid the exponential objective's worst-case gradient growth. The loss is selectable via \texttt{LOSS\_FN=asymmetric} or \texttt{--loss-fn asymmetric} in \texttt{client.py}, \texttt{final\_benchmark.py}, and \texttt{fed\_sim.py}; the default remains \texttt{mse}. For prediction residual $d = \hat{y} - y$, it applies a $C^1$-smooth asymmetric quadratic penalty:
\begin{equation}
\mathcal{L}_{\text{asym}}(d) = \begin{cases}
\alpha_{\text{late}}\,d^2, & d \ge 0 \quad (\text{late prediction: }\hat{y} \ge y),\\
\alpha_{\text{early}}\,d^2, & d < 0 \quad (\text{early prediction: }\hat{y} < y),
\end{cases}
\end{equation}
with $\alpha_{\text{late}} > \alpha_{\text{early}} > 0$ (default $\alpha_{\text{late}}{=}2.5$, $\alpha_{\text{early}}{=}0.5$, a $5\times$ penalty ratio on late predictions). Because $\mathcal{L}_{\text{asym}}$ is quadratic on both half-lines, its gradient is continuous at $d=0$ and bounded linearly in $|d|$, ensuring stable Adam updates while steering both local and federated models away from dangerous RUL overestimation.

\subsection{Evaluation metrics}
With $\Delta_i=\hat y_i-y_i$ (positive = late prediction), RMSE is $\sqrt{\frac1M\sum_i\Delta_i^2}$, the late-prediction rate (\texttt{late\_prediction\_rate} in \texttt{common.py}) is $\frac{1}{M}\sum_{i=1}^{M}\mathbb{I}(\hat{y}_i > y_i)$, and the NASA score~\cite{saxena2008damage} is
\begin{equation}
S=\sum_{i=1}^{M}\begin{cases}e^{-\Delta_i/13}-1,&\Delta_i<0\\ e^{\Delta_i/10}-1,&\Delta_i\ge 0.\end{cases}
\end{equation}
Late errors are weighted by the smaller constant 10 and so grow faster. $S$ is a sum and thus is not comparable across test sets of different size (FD001 vs.\ FD002).

\subsection{Local degradation attribution and healthy-baseline drift}
\label{sec:method-diag}
\textbf{Channel map.} The model input has 17 channels: three operating settings and 14 sensors. Following Table~2 of Saxena et al.~\cite{saxena2008damage}, the 14 sensors are mapped to four subsystems (Table~\ref{tab:channels}). The settings \texttt{setting1..3} remain model inputs but describe the flight condition, not the engine's health, so they are explicitly excluded from fault responsibility. The Turbine group includes the fuel-flow ratio $\phi$ (s12), which is defined relative to the HPC outlet pressure; we place it with the combustor and hot section. This assignment affects the HPC-versus-Turbine split reported later, but not the core-versus-Fan/LPC split.

\begin{table}[t]
\centering
\caption{Sensor-to-subsystem map used for attribution, after Table~2 of~\cite{saxena2008damage}. Operating settings carry no fault responsibility.}
\label{tab:channels}
\footnotesize
\begin{tabular}{@{}lp{0.74\columnwidth}@{}}
\toprule
Subsystem & Sensors \\
\midrule
Fan & Nf (s8), NRf (s13), BPR (s15) \\
LPC & T24 (s2) \\
HPC & T30 (s3), P30 (s7), Nc (s9), Ps30 (s11), NRc (s14) \\
Turbine & T50 (s4), $\phi$ (s12), htBleed (s17), W31 (s20), W32 (s21) \\
Excluded & \texttt{setting1}, \texttt{setting2}, \texttt{setting3} \\
\bottomrule
\end{tabular}
\end{table}

\textbf{Baseline-relative attribution.} For a window $\mathbf{x}\in\mathbb{R}^{W\times 17}$ in the model's scaled input space and its prediction $\hat y=f_{\mathbf{w}}(\mathbf{x})$, the attribution of channel $c$ is
\begin{equation}
A_c=\frac{1}{W}\sum_{t=1}^{W}\left|\bigl(x_{t,c}-b_{t,c}\bigr)\,\frac{\partial \hat y}{\partial x_{t,c}}\right|,
\label{eq:attr}
\end{equation}
where $\mathbf{b}$ is the engine's healthy-baseline expectation for the same cycles and operating regime (defined below). This is a first-order estimate of how much the deviation of channel $c$ from healthy moves the RUL estimate. It is a baseline variant of gradient$\times$input attribution~\cite{simonyan2014saliency,shrikumar2017deeplift,sundararajan2017axiomatic}; plain saliency $|\partial\hat y/\partial x_{t,c}|$ is kept as a comparison. For the settings, $\mathbf{b}$ equals the observed values, and their $A_c$ is additionally set to zero. Sensor shares are $s_c=100\,A_c/\sum_{c'}A_{c'}$ over the 14 sensors, and a subsystem share is the sum of its sensors' shares. Both therefore sum to exactly 100\%. If every $A_c$ is zero, the shares fall back to a uniform split. Displayed one-decimal shares use largest-remainder rounding, so they still sum to exactly 100.0. The \emph{primary} subsystem of a window is the one with the largest share.

\textbf{Regime-aware healthy-baseline drift.} The drift score does not use the model. Each raw cycle is assigned to the nearest of $k$ operating-regime centroids, fitted by k-means on standardized settings ($k{=}1$ for FD001/FD003, $k{=}6$ for FD002/FD004). Each sensor is standardized with that regime's mean and standard deviation, giving $u_{t,c}$. The first $H{=}30$ cycles of each engine define its healthy offset, and
\begin{equation}
z_{t,c}=\frac{u_{t,c}-\frac{1}{H}\sum_{t'\le H}u_{t',c}}{\sigma^{\text{h}}_c},
\label{eq:drift}
\end{equation}
where $\sigma^{\text{h}}_c$ is the pooled within-engine early-life noise of the local fleet. Window drift is the mean of $z_{t,c}$ over the window. The baseline $\mathbf{b}$ in~\eqref{eq:attr} is the same regime mean plus engine offset, mapped back to raw units and scaled like the model input. The subsystem drift is the mean $|z|$ of its sensors, which does not depend on how many sensors a subsystem has.

\textbf{Maintenance summary.} For each window, the module reports the predicted RUL, a severity level from fixed predicted-RUL thresholds, the subsystem shares, the top-3 contributing and top-3 drifting sensors with their physical names, and the evaluation-boundary note that the output is an explanation, not a supervised diagnosis.
